\par\noindent\textbf{(1)} 远场复振幅与孔面均匀透过函数的二维Fourier变换成正比。令横向波矢 \(q=k\sin\theta\)，角积分 \(\int_0^{2\pi}e^{iqr\cos\phi}d\phi=2\pi J_0(qr)\)，故振幅 \(A(q)\propto2\pi\int_0^arJ_0(qr)dr\)。
\par\noindent\textbf{(2)} 使用 \(\int_0^arJ_0(qr)dr=aJ_1(qa)/q\)，归一复振幅 \(A(\theta)/A(0)=2J_1(x)/x\)，其中 \(x=ka\sin\theta\)。故强度 \(I(\theta)/I(0)=[2J_1(x)/x]^2\)，中心按 \(J_1(x)\sim x/2\) 连续趋1。
\par\noindent\textbf{(3)} Bessel函数 \(J_1\) 首个正零点为 \(x_1\approx3.831706\)，因此 \(\sin\theta_1=x_1/(ka)\approx1.22\lambda/D\)，\(D=2a\)。Rayleigh判据将一个点像中心放置到另一个点像第一暗环作为分辨参考尺度，但不是严格的信息理论极限。
\par\noindent\textbf{(4)} 当 \(D\) 减半，第一暗环角约增加一倍，中心主瓣扩展。透过孔径的总功率按面积 \(\pi a^2\) 下降，但归一化衍射强度图形仍服从 \(x=ka\sin\theta\) 的自相似尺度。不可将远场中心归一强度保持1误解释为总能量不变。
